% Figure: virtual work on a simple lever / two-bar linkage. A virtual
% displacement of the input rotates the mechanism; the principle of
% virtual work balances input and output through the geometry alone, with
% no need to compute the pin reaction.
\begin{figure}[htbp]
\centering
\begin{tikzpicture}[
  bar/.style={draw=engblack, very thick},
  ghost/.style={draw=enggray, very thick, dashed},
  force/.style={draw=engred, -{Stealth}, very thick},
  pin/.style={draw=engblack, fill=white, circle, inner sep=1.6pt},
  vd/.style={draw=engblue, -{Stealth}, thick},
  label/.style={font=\small},
]
% Fulcrum at origin
\node[pin] (F) at (0,0) {};
\node[font=\scriptsize, below] at (0,-0.15) {fulcrum};

% Lever: left arm length a = 3.4 (input), right arm length b = 1.7 (output)
% horizontal at rest
\coordinate (A) at (-3.4, 0);
\coordinate (B) at (1.7, 0);
\draw[bar] (A) -- (B);

% Virtual rotation delta-phi: ends move vertically by a*dphi and b*dphi
% rotate by small angle (drawn at 12 deg for visibility)
\coordinate (Ap) at ({-3.4*cos(12)}, {-3.4*sin(12)});
\coordinate (Bp) at ({1.7*cos(12)}, {1.7*sin(12)});
\draw[ghost] (Ap) -- (Bp);

% input force F1 downward at A (load pushes A down through delta r1 = a dphi)
\draw[force] (-3.4, 1.3) -- (-3.4, 0.15);
\node[engred, font=\small, anchor=south] at (-3.4, 1.35) {$F_1$};
% output force F2 downward at B resisting
\draw[force] (1.7, -1.3) -- (1.7, -0.15);
\node[engred, font=\small, anchor=north] at (1.7, -1.35) {$F_2$};

% virtual displacements
\draw[vd] (-3.4, 0) -- (Ap);
\node[engblue, font=\scriptsize, anchor=east] at ($(Ap)+(-0.05,0)$) {$\delta r_1 = a\,\delta\phi$};
\draw[vd] (1.7, 0) -- (Bp);
\node[engblue, font=\scriptsize, anchor=west] at ($(Bp)+(0.05,0)$) {$\delta r_2 = b\,\delta\phi$};

% arc showing delta phi at fulcrum
\draw[draw=englime, -{Stealth}, thick] (0.9,0) arc (0:12:0.9);
\node[englime, font=\scriptsize, anchor=west] at (0.95, 0.15) {$\delta\phi$};

\end{tikzpicture}
\caption[Principle of virtual work on a lever]{The principle of virtual
work applied to a rigid lever pivoted at a frictionless fulcrum. A
virtual rotation $\delta\phi$ (green) consistent with the pin constraint
moves the left end down by $\delta r_1=a\,\delta\phi$ and the right end
up by $\delta r_2=b\,\delta\phi$ (blue). The pin reaction does no virtual
work because the fulcrum does not move. Setting the total virtual work
of the applied forces to zero, $F_1\,a\,\delta\phi - F_2\,b\,\delta\phi=0$,
gives the equilibrium law $F_1 a = F_2 b$ directly, without ever solving
for the reaction at the pin. Static equilibrium is the zero-velocity
limit of the energy method: the workless constraints drop out, leaving a
scalar balance read off the geometry.}
\label{fig:vol03ch05:virtual-work}
\end{figure}
